\scanpage{156}

\[
\det(t\sigma_\infty+u)=t^2+u^2=\varepsilon_2/\mu
\quad(=1/\mu\text{ if }\varepsilon_2=+1),
\]
we shall indeed have, in the case $\varepsilon_2=1$, \deleted{\unsure{therefore}}
\[
\det\uline{u}=\mu^2\det(t\sigma_\infty+u)=\mu^2\,1/\mu=\mu
\]
i.e., (66). On the other hand, we shall have
\[
\begin{aligned}
[u,\rho]&=[t\sigma_\infty+u,\rho]\\
&=
\underbrace{
\begin{pmatrix}u&t\\-t&u\end{pmatrix}
\begin{pmatrix}0&1\\-1&1\end{pmatrix}
}_{\begin{pmatrix}-t&u+t\\-u&u-t\end{pmatrix}}
\inserted{\mu}
\underbrace{
\begin{pmatrix}u&-t\\t&u\end{pmatrix}
\begin{pmatrix}1&-1\\1&0\end{pmatrix}
}_{\begin{pmatrix}u-t&-u\\u+t&-t\end{pmatrix}}.
\\
&=\mu\begin{pmatrix}
(u+t)^2-ut+t^2&-t^2\\
ut-t^2&u^2+t^2-ut
\end{pmatrix}
=\mu\begin{pmatrix}
\dfrac1\mu+\unsure{t(u+t)}&-t^2\\
t(u-t)&\dfrac1\mu-ut
\end{pmatrix}
\\
&=\begin{pmatrix}
1+\mu t(u+t)&-\mu t^2\\
\mu t(u\deleted{\unsure{t}}-t^{\deleted{\unsure{2}}})&1-\mu ut
\end{pmatrix},
\end{aligned}
\]
\margin{We could also \unsure{apply} the formulas §43 (78) ---}
and comparing with (65) we do indeed find
\[
[u,\rho]=g
\]
taking into account that
\[
\mu(u^2+t^2)=\varepsilon_2=1
\]
thus
\[
-1+\mu u(t+u)=\underbrace{-1+\mu u^2}_{-\mu t^2}+\mu ut
=\mu t(u-t)\ \unsure{-!|!}
\]
I would now like, without launching into
new calculations, to find ``by guesswork'' an expression
for $u$, in the \deleted{\unsure{case}} where $\varepsilon_2=-1$, and where we expect
to find in (67) the sign $\varepsilon'=-1$ --- i.e. that
$u$ be an \uline{anti}\unsure{similitude}\illegible{1}
\[
\uline{u}\,{}^t\uline{u}=-\det\uline{u}
\]
But \deleted{\unsure{and}} the matrix
\[
\tau_\infty=\begin{pmatrix}-1&0\\0&1\end{pmatrix}
\]
\margin{(it is \unsure{$=$} an \unsure{orthogonal} reflection)}
is the \unsure{typical case} of such an antisimilitude, and I
expect that the matrix
\[
\uline{u}=\mu\tau_\infty(t\sigma_\infty+u)
=\mu\begin{pmatrix}\unsure{1|-1}&0\\0&-1\end{pmatrix}
\begin{pmatrix}u&t\\-t&u\end{pmatrix}
=\mu\begin{pmatrix}\unsure{u|-u}&\unsure{t|-t}\\-t&-u\end{pmatrix}
\]
\margin{155}

\scanpage{157}
\mspage{368}

will do. We shall have, of course
\[
\begin{aligned}
u\,{}^tu
&=\mu^2\begin{pmatrix}u&t\\t&-u\end{pmatrix}
\begin{pmatrix}u&t\\t&-u\end{pmatrix}\\
&=\mu^2\begin{pmatrix}u^2+t^2&0\\0&u^2+t^2\end{pmatrix}
=\mu^2\begin{pmatrix}-1/\mu&0\\0&-1/\mu\end{pmatrix}
\end{aligned}
\]
i.e.
\[
u\,{}^tu=-\mu.\mathrm{Id}
\]
and
\[
\det u=\mu^2(-u^2-t^2)=\mu^2(1/\mu)=\mu
\]
taking into account that now
\[
u^2+t^2=\varepsilon_2/\mu=-\frac1\mu
\qquad\text{since }\varepsilon_2=-1.
\]
Finally to check (68) let us compute
\[
\begin{aligned}
[u,\rho]
&=\left[\begin{pmatrix}\unsure{u}&t\\t&-u\end{pmatrix},\rho\right]\\
&=\underbrace{
\begin{pmatrix}u&t\\t&-u\end{pmatrix}
\begin{pmatrix}0&1\\-1&1\end{pmatrix}
}_{\begin{pmatrix}-t&u+t\\u&-u+t\end{pmatrix}}
\unsure{\dfrac1\mu|\mu}
\underbrace{
\begin{pmatrix}-u&-t\\-t&u\end{pmatrix}
\begin{pmatrix}1&-1\\1&0\end{pmatrix}
}_{\begin{pmatrix}-u-t&\unsure{u|-u}\\u-t&t\end{pmatrix}}
\\
&=\mu\begin{pmatrix}
u(u+t)&t^2\\
\unsure{-(u^2+t^2)}-u(u-t)&u^2+t^2-ut
\end{pmatrix}
=\begin{pmatrix}
\mu u(u+t)&\mu t^2\\
1-\mu u(u-t)&-1-\mu ut
\end{pmatrix}
\end{aligned}
\]
(taking into account that $\mu(u^2+t^2)=-1$), \deleted{\unsure{whence|therefore}} do we indeed have
\[
[u,\rho]=-g
\]
taking
into account formula (65) which gives $g$ explicitly in
terms of $r,s,t,u$, we find that this
\unsure{expected|precise} formula is equivalent to
\[
2\mu ut=0\qquad\text{i.e.}\qquad ut=0
\]
really, this is screwy! But I see that if we
replace in the definition of $u$ $t$ by $-t$,
i.e. if we define
\[
\uline{u}=\mu(t\sigma_\infty+u)\tau_\infty
=\mu\begin{pmatrix}u&t\\-t&u\end{pmatrix}
\begin{pmatrix}1&0\\0&-1\end{pmatrix}
=\mu\begin{pmatrix}u&-t\\-t&-u\end{pmatrix}
\]
(instead of $\uline{u}=\mu\begin{pmatrix}u&t\\t&-u\end{pmatrix}$),
then $\uline{u}$ is \unsure{always}
an antisimilitude of determinant $\mu$, and
moreover now
\[
[u,\rho]=\begin{pmatrix}
\mu u(u-t)&\mu t^2\\
1-\mu u(u+t)&-1+\mu ut
\end{pmatrix}=-g
\]
\margin{156}

\scanpage{158}

which does indeed reduce to checking
\[
\mu u(u-t)=-(1+\mu t(t+u))
\]
i.e.
\[
\mu[u^2-\deleted{ut}+t^2+\deleted{tu}]=-1
\qquad\text{O.K.}
\]
In summary, we find

\uline{Theorem} Let $\alpha_2\in\mathbb Z/2\mathbb Z$ be defined by
$(-1)^{\alpha_2}=\varepsilon_2$
(thus $\alpha_2\equiv0\ (2)$ iff $\mu\equiv1\ (4)$,
$\alpha_2\equiv1\ (2)$ iff
$\mu\equiv-1\ (4)$). Put
\[
\uline{u}=\mu(t\sigma_\infty+u)\tau_\infty^{\alpha_2}
=\mu\begin{pmatrix}u&t\\-t&u\end{pmatrix}
\begin{pmatrix}1&0\\0&\varepsilon_2\end{pmatrix}^{\unsure{\alpha_2}},
\tag{75}
\]
\margin{$\displaystyle
=\begin{pmatrix}u&\varepsilon_2t\\-t&\varepsilon_2u\end{pmatrix}$}
\margin{\deleted{$\displaystyle
=\left\{\begin{array}{ll}
\mu\begin{pmatrix}u&t\\-t&u\end{pmatrix}&\alpha_2=0\\
\mu\begin{pmatrix}u&-t\\-t&-u\end{pmatrix}&\alpha_2=1
\end{array}\right.$}}
We then have the formulas
\[
\det\uline{u}=\mu
\tag{76}
\]
\[
[\uline{u},\sigma_\infty]=\unsure{-|\cdot}\varepsilon_2.\mathrm{id}
\qquad\text{i.e.}\quad\uline{u}(\sigma_\infty)=\varepsilon_2.\sigma_\infty
\tag{77}
\]
\[
[\uline{u},\rho]=\varepsilon_2.g
\qquad\text{i.e.}\quad\uline{u}(\rho)=(\varepsilon_2g).\rho
\tag{78}
\]

\uline{Corollary} Let $u$ be an automorphism \unsure{of the} loop
group $\widehat{\Pi}_{0,3}$, \unsure{commuting} as outer actions with
$\unsure{\mathfrak S_3|\mathcal S_3}$ (i.e. with $\sigma_\infty$ and with $\rho$), $\mu$
its multiplier, let us also denote by
$u$ the canonical extension of $u$ to
$\unsure{\widehat{\mathfrak S}_{0,3}^{+}|\widehat{\mathcal S}_{0,3}^{+}}
\simeq\mathrm{Sl}(2,\mathbb Z)^{\wedge\prime}$,
inducing the identity on
$\unsure{\mathfrak S_3|\mathcal S_3}
\simeq\unsure{\widehat{\mathfrak S}_{0,3}^{+}|\widehat{\mathcal S}_{0,3}^{+}}/
\widehat{\Pi}_{0,3}$,
and consider the quotient $\mathrm{Sl}(2,\widehat{\mathbb Z})'$ of
$\mathrm{Sl}(2,\mathbb Z)^{\wedge\prime}$. Then the action of $u$ passes
to an action on the quotient $\mathrm{Sl}(2,\widehat{\mathbb Z})'$,
and the latter is induced by an
\deleted{\unsure{inner}}\inserted{``\unsure{inner|inner}''}
automorphism of $\mathrm{Gl}(2,\widehat{\mathbb Z})$, defined by a
\inserted{matrix} $\uline{u}\in\mathrm{Gl}(2,\widehat{\mathbb Z})$
such that $\det u=\mu$
\margin{157}

\scanpage{159}
\mspage{369}

(which describes the \uline{outer} action of $u$ on
$\mathrm{Sl}(2,\widehat{\mathbb Z})'$ in terms of the multiplier $\mu$).

\margin{\uline{N.B.} To check the \unsure{corollary}, we are
reduced to the case where $u$ commutes with $\sigma_\infty$,
if necessary conjugating by an inner automorphism \unsure{of}\illegible{1}\ldots}
\margin{We have done this, for the moment, for $u$
commuting with $\sigma_\infty$.}

In fact, we have much better still,
since we are going to define \uline{canonically}
\inserted{\uline{modulo sign}},
in terms of $u$, the matrix $\uline{u}$ by
formula (75), where \deleted{$t\text{ and }u\in$} $t,u\in\widehat{\mathbb Z}$ are
defined as follows: \deleted{\unsure{We have}} the datum of the automorphism $u$
defines $f\in\widehat{\Pi}_{0,3}$, modulo the operation $f\mapsto\ell_0^n f$
(for $n\in\widehat{\mathbb Z}$), and we again denote by $f$
the image of $f$ in $\mathrm{Sl}(2,\widehat{\mathbb Z})/\pm1$, which is thus
a matrix $\begin{pmatrix}r&s\\t&u\end{pmatrix}$, defined modulo sign,
We shall note that in $\mathrm{Sl}(2,\widehat{\mathbb Z})$ we have
\[
\ell_0^n\begin{pmatrix}r&s\\t&u\end{pmatrix}
=\begin{pmatrix}1&-2n\\0&1\end{pmatrix}
\begin{pmatrix}r&s\\t&u\end{pmatrix}
=\begin{pmatrix}r-2nt&s-2nu\\t&u\end{pmatrix}
\tag{79}
\]
($n\in\widehat{\mathbb Z}$), so $(t,u)$ does not change (modulo sign)
when we replace $f$ by $\ell_0^n f$ --- so $(t,u)$
is determined modulo sign by the automorphism $u$
of $\widehat{\Pi}_{0,3}$. \deleted{\illegible{1}} The analogous observation applies
moreover to $\alpha_0=\begin{pmatrix}a&b\\c&d\end{pmatrix}$, which is determined
modulo sign and modulo left multiplication
by an $\ell_0^n$ --- the \deleted{\illegible{1}} vector $(c,d)$ is
determined modulo sign, hence \unsure{the element}
\[
c\rho+d
\]
\margin{158}

\scanpage{160}

is determined modulo \uline{sign}, and \deleted{\illegible{1}}\inserted{unless I am mistaken}
we must have a relation
\[
\tag{80}
\uline{u}\overset{?}{=}\mu\alpha(c\rho+d)
=\mu f^{-1}\ell_0^{-\nu}\alpha_0(c\rho+d)
\]
(cf. (69) (73)), at least if $\varepsilon_2=1$ ---
otherwise we may perhaps have to correct by a factor $\tau_\infty$ somewhere.
\deleted{Now I have no \unsure{desire} for the moment to resume \unsure{my|my} \unsure{calculations|calculation} in this direction.}

\margin{N.B. this formula is \unsure{valid} \unsure{with} the corrective factor $\delta=\varepsilon_3$, cf. (82) below\ldots}

\[
\ell_0^{-\nu}\alpha_0(c\rho+d)
=
\begin{pmatrix}1&2\nu=\mu-1\\0&1\end{pmatrix}
\underbrace{
\begin{pmatrix}a&b\\c&d\end{pmatrix}
\begin{pmatrix}d&c\deleted{\unsure{+d}}\\-c&c+d\end{pmatrix}
}_{\displaystyle
\underbrace{
\begin{pmatrix}
ad-bc=\delta&ac+bc+bd=\delta B\\
0&c^2+cd+d^2=\delta D
\end{pmatrix}
}_{\displaystyle\delta\begin{pmatrix}1&B\\0&D\end{pmatrix}}
}
\]
\[
\ell_0^{-\nu}\alpha_0(c\rho+d)
=\delta\begin{pmatrix}1&B+2\nu D=C\\0&D\end{pmatrix}
\]
i.e.
\[
\tag{81}
\ell_0^{-\nu}\alpha_0(c\rho+d)
=\delta\begin{pmatrix}1&C\\0&D\end{pmatrix}
\]
whence
\[
\begin{aligned}
\mu\alpha(c\rho+d)
&=(\mu f^{-1})(\ell_0^{-\nu}\alpha_0(c\rho+d))
=\mu\delta\begin{pmatrix}u&-s\\-t&r\end{pmatrix}
\begin{pmatrix}1&C\\0&D\end{pmatrix}\\
&=\mu\delta\begin{pmatrix}u&uC-sD\\-t&-tC+rD\end{pmatrix}\\
&=\delta\begin{pmatrix}\mu u&-s+\mu uC\\-\mu t&r-\mu tC\end{pmatrix}
\end{aligned}
\]
Now (61)
\[
C=\varepsilon_2T=\varepsilon_2(rt+su)\qquad\text{whence}
\]
\[
\begin{aligned}
-s+\mu uC
&=-s+\varepsilon_2\unsure{\mu urt|\mu turt}+\varepsilon_2\mu u^2s\\
&\overset{\inserted{\unsure{\text{if }\varepsilon_2=+1}}}{=}
 s\bigl(\underbrace{-1+\unsure{\varepsilon_2u^2|\varepsilon_2\mu u^2}}_{-\varepsilon_2\mu t^2}\bigr)
 +\varepsilon_2\mu urt\\
&=\varepsilon_2\mu t\bigl(\overbrace{\unsure{-st+ur|st+ur}}^{+1}\bigr)\\
&=\varepsilon_2\mu t
\end{aligned}
\]

\margin{159}

\scanpage{161}
\mspage{370}

\[
\begin{aligned}
r-\mu tC
&\overset{\inserted{\unsure{\text{if }\varepsilon_2=1}}}{=}
\underbrace{r-\varepsilon_2\mu rt^2}_{r(1-\varepsilon_2\mu t^2)=\varepsilon_2r\mu u^2}
-\varepsilon_2\mu stu\\
&=\varepsilon_2\mu u(ru-st)=\varepsilon_2\mu u,
\end{aligned}
\]
thus
\[
\mu\alpha(c\rho+d)
=\delta\mu\begin{pmatrix}u&\varepsilon_2\unsure{t}\\-t&\varepsilon_2u\end{pmatrix}
\overset{\deleted{\unsure{\text{if }\varepsilon_2=1}}}{=}\delta\uline{u},
\]
i.e.
\[
\tag{82}
\uline{u}=(\unsure{\delta}\mu)\alpha(c\rho+d)
=(\delta\mu)(f^{-1}\ell_0^{-\nu}\alpha_0)(c\rho+d),
\]
\deleted{\unsure{if $\varepsilon_2=1$ (i.e. $\mu\equiv1\ (4)$)}}
\inserted{where $\delta=ad-bc\in\{\pm1\}$, $\delta=\varepsilon_3$
(thus $\delta\equiv\mu\ (3)$)}.

Formula which can also be written
\deleted{i.e.} with \unsure{account} \deleted{\illegible{1}} taken of (75)
\[
\tag{83}
\left\{\begin{aligned}
\deleted{\unsure{f}}(t\sigma_\infty+u)\inserted{\tau_\infty^{\alpha_2}}
&=\delta\alpha(c\rho+d)
\qquad\deleted{\unsure{\text{if }\varepsilon_2=1}}\\
\text{where}\quad\alpha&=f^{-1}\ell_0^{-\nu}\alpha_0,
\quad\deleted{\illegible{1}}\quad
\delta=\varepsilon_3\in\{\pm1\},\quad\delta=ad-bc,\quad\delta\equiv\mu\ (3)
\end{aligned}\right.
\]
or again
\[
\tag{84}
f(t\sigma_\infty+u)\tau_\infty^{\alpha_2}
=\delta\ell_0^{-\nu}\alpha_0(c\rho+d)
\]
\[
\begin{array}{c|c}
\underbrace{
\begin{pmatrix}r&s\\t&u\end{pmatrix}
\begin{pmatrix}u&t\\-t&u\end{pmatrix}
}_{\displaystyle\begin{pmatrix}1&rt+su\\0&t^2+u^2\end{pmatrix}}
\begin{pmatrix}1&0\\0&\varepsilon_2\end{pmatrix}
&
\underbrace{
\begin{pmatrix}1&2\nu\\0&1\end{pmatrix}
\begin{pmatrix}a&b\\c&\unsure{d}\end{pmatrix}
\begin{pmatrix}d&c\\-c&c+d\end{pmatrix}
}_{\displaystyle\begin{pmatrix}1&C\\0&D\end{pmatrix}}
\quad\text{cf. (81)}\\
\begin{pmatrix}1&rt+su\\0&t^2+u^2\end{pmatrix}\deleted{\begin{pmatrix}\unsure{u}&\unsure{t}\\-\unsure{t}&\unsure{u}\end{pmatrix}}&\\
\deleted{=\begin{pmatrix}1&T\\0&U\end{pmatrix}}&\\
\deleted{\unsure{=}}&\\
=\begin{pmatrix}1&\varepsilon_2T\\0&\varepsilon_2U\end{pmatrix}&
\end{array}
\]
and in this form the formula \unsure{was} in fact known,
as the essential content of the \deleted{formula} loop relation in matrix form (56) or (57)
\[
:\quad\varepsilon_2U=D,\qquad\varepsilon_2T=C\qquad!
\]

\uline{Remark} In fact, under a false appearance of arithmetic, we have done \unsure{by|for} the back door some scales in linear algebra, valid over

\margin{160}

\scanpage{162}

arbitrary rings and \unsure{formulable|formulable} probably under much more general and more ``\unsure{comprehensive}'' conditions --- concerning equations of the form
\[
[\alpha_0,\rho]=\xi
\]
where $\rho\in\unsure{\mathrm{Gl}}(2,A)$ is given, \deleted{\unsure{and}\illegible{1}}\inserted{as well as} $\xi$, and
$\alpha_0\in\mathrm{Gl}(2,A)$ is the unknown (possibly subject to conditions of the type $\det\alpha_0=\delta$, $\delta$ fixed), as well as simultaneous equations in $\alpha_0$ and $f$
\[
\tag{*}
[\alpha_0,\rho]=\varepsilon.\unsure{\mathcal{E}_0|\varepsilon_0}(\mu-1)[f,\sigma]
\qquad\text{cf. (56)}
\]
where $\rho,\sigma\in\mathrm{Gl}(2,A)$ are given, as well as a scalar $\varepsilon\in A^*$, and
$\lambda\longmapsto\unsure{\mathcal{E}_0|\varepsilon_0}(\lambda)$ is a ``one-parameter subgroup'', for example we can assume
\[
\unsure{\mathcal{E}_0|\varepsilon_0}(\lambda)=\begin{pmatrix}1&-\lambda\\0&1\end{pmatrix}=\varepsilon_0^\lambda,
\]
\[
\varepsilon_0=\begin{pmatrix}1&-1\\0&1\end{pmatrix}
\]
--- assuming if necessary a relation of the form $\sigma\rho=\varepsilon_0$.

Modulo some pleasant adjustments, all the formulas of the present \S\ should be able to be developed in this more general context
\inserted{(as well as substantial complements)}
(--- but I do not know whether I am going to spell out

\margin{161}